\section{Introduction}

Neural networks are typically modeled as statistical correlation machines. We argue they are better understood as physical systems seeking homeostatic equilibrium.

In this work, we synthesize evidence from three foundational studies to present a unified theory of \emph{Homeostatic Phase Transitions}.
\begin{enumerate}
    \item \textbf{Anatomy:} \cite{thompson2025layer0} established that "Layer-0 Suppressors" act as the braking system, trading factuality for hedging in large models like GPT-2 and Mistral.
    \item \textbf{Dynamics:} \cite{thompson2025compensation} demonstrated "Active Homeostatic Compensation," showing that networks obey Le Chatelier's Principle—actively redistributing variance to oppose perturbations.
    \item \textbf{Limits:} \cite{thompson2025timescale} mapped the "Observable-Dependent Landscapes," proving that these dynamics are constrained by information saturation boundaries.
\end{enumerate}

Building on these pillars, we report a new, unifying discovery: these diverse phenomena collapse onto a single, precise geometric reality.

\paragraph{The 12-Decimal Anomaly.}
We discover that when trained on algorithmic tasks, transformer attention converges to a conserved quantity independent of initialization. Five independent seeds converged to an Effective Rank index of exactly \textbf{0.611991643906}—precise to 12 decimal places.

This is not a statistical average; it is a chaotic attractor collapsing to a fixed point. In this paper, we identify the mechanism of this collapse (Section~\ref{sec:mechanism}), map its geometric limits (Section~\ref{sec:geometry}), and exploit it to accelerate training by 93\% (Section~\ref{sec:engineering}).
